\subsection{Independent principal-angle calculation for the arithmetic four-volume}\label{endpoint-principal_angle_volume_review--independent-principal-angle-calculation-for-the-arithmetic-four-volume}

13 September 2026. This is a separate written derivation of the proposed canonical-section angle formula and its integrated metric comparison. No sealed source is edited.

\subsubsection{1. Canonical sections in one original Hilbert space}\label{endpoint-principal_angle_volume_review--canonical-sections-in-one-original-hilbert-space}

Fix the original monic polynomial ENDPOINTSOURCE0C01A7DD4A656A73579A2F43SLOT0END of degree ENDPOINTSOURCE0C01A7DD4A656A73579A2F43SLOT1END, the original real weight ENDPOINTSOURCE0C01A7DD4A656A73579A2F43SLOT2END, and the literal finite polynomial space

ENDPOINTSOURCE0C01A7DD4A656A73579A2F43SLOT3END

with its fixed coefficient basis and remainder map. Let ENDPOINTSOURCE0C01A7DD4A656A73579A2F43SLOT4END be a continuously differentiable Gram on this common space. All orthogonal projections and norms below use that Gram at the same parameter value. The adjoint in the coefficient formula ENDPOINTSOURCE0C01A7DD4A656A73579A2F43SLOT5END is the ordinary conjugate transpose in the retained bases; no coefficient basis is changed implicitly.

For ENDPOINTSOURCE0C01A7DD4A656A73579A2F43SLOT6END, let ENDPOINTSOURCE0C01A7DD4A656A73579A2F43SLOT7END project onto ENDPOINTSOURCE0C01A7DD4A656A73579A2F43SLOT8END, and let ENDPOINTSOURCE0C01A7DD4A656A73579A2F43SLOT9END project onto the original relation space ENDPOINTSOURCE0C01A7DD4A656A73579A2F43SLOT10END, with ENDPOINTSOURCE0C01A7DD4A656A73579A2F43SLOT11END. Since the latter is a subspace of the former, their projections commute and satisfy ENDPOINTSOURCE0C01A7DD4A656A73579A2F43SLOT12END. Therefore

ENDPOINTSOURCE0C01A7DD4A656A73579A2F43SLOT13END

is exactly the orthogonal projection onto ENDPOINTSOURCE0C01A7DD4A656A73579A2F43SLOT14END. The original least-norm remainder section ENDPOINTSOURCE0C01A7DD4A656A73579A2F43SLOT15END has ENDPOINTSOURCE0C01A7DD4A656A73579A2F43SLOT16END, image ENDPOINTSOURCE0C01A7DD4A656A73579A2F43SLOT17END, and positive Gram

ENDPOINTSOURCE0C01A7DD4A656A73579A2F43SLOT18END

These assertions follow from the direct orthogonal decomposition of the surjective remainder sequence; its kernel is multiplication by the unchanged monic ENDPOINTSOURCE0C01A7DD4A656A73579A2F43SLOT19END.

If ENDPOINTSOURCE0C01A7DD4A656A73579A2F43SLOT20END, both ENDPOINTSOURCE0C01A7DD4A656A73579A2F43SLOT21END and ENDPOINTSOURCE0C01A7DD4A656A73579A2F43SLOT22END are in ENDPOINTSOURCE0C01A7DD4A656A73579A2F43SLOT23END and have remainder ENDPOINTSOURCE0C01A7DD4A656A73579A2F43SLOT24END. Their difference belongs to ENDPOINTSOURCE0C01A7DD4A656A73579A2F43SLOT25END, which is orthogonal to ENDPOINTSOURCE0C01A7DD4A656A73579A2F43SLOT26END. Hence

ENDPOINTSOURCE0C01A7DD4A656A73579A2F43SLOT27END

The second identity follows by writing ENDPOINTSOURCE0C01A7DD4A656A73579A2F43SLOT28END; its last summand is orthogonal to ENDPOINTSOURCE0C01A7DD4A656A73579A2F43SLOT29END. The same expansion gives the last identity with all cross terms accounted for.

\subsubsection{2. Exact principal angles and the square-root order}\label{endpoint-principal_angle_volume_review--exact-principal-angles-and-the-square-root-order}

Define the actual isometries

ENDPOINTSOURCE0C01A7DD4A656A73579A2F43SLOT30END

Their domains have the usual coefficient inner product; their targets carry ENDPOINTSOURCE0C01A7DD4A656A73579A2F43SLOT31END. Thus ENDPOINTSOURCE0C01A7DD4A656A73579A2F43SLOT32END. Their cross-Gram and its two products, in their exact noncommutative orders, are

ENDPOINTSOURCE0C01A7DD4A656A73579A2F43SLOT33END

The first follows from (PA4), with ENDPOINTSOURCE0C01A7DD4A656A73579A2F43SLOT34END in the adjacent factors. No commutation between ENDPOINTSOURCE0C01A7DD4A656A73579A2F43SLOT35END and ENDPOINTSOURCE0C01A7DD4A656A73579A2F43SLOT36END is used.

The singular values of the first matrix in (PA6) are the cosines of the principal angles between ENDPOINTSOURCE0C01A7DD4A656A73579A2F43SLOT37END and ENDPOINTSOURCE0C01A7DD4A656A73579A2F43SLOT38END. This can be seen directly. Choose an orthonormal eigenbasis ENDPOINTSOURCE0C01A7DD4A656A73579A2F43SLOT39END of ENDPOINTSOURCE0C01A7DD4A656A73579A2F43SLOT40END, and let its eigenvalues be ENDPOINTSOURCE0C01A7DD4A656A73579A2F43SLOT41END. By (PA3)--(PA4), ENDPOINTSOURCE0C01A7DD4A656A73579A2F43SLOT42END. Put

ENDPOINTSOURCE0C01A7DD4A656A73579A2F43SLOT43END

The ENDPOINTSOURCE0C01A7DD4A656A73579A2F43SLOT44END form an orthonormal basis of ENDPOINTSOURCE0C01A7DD4A656A73579A2F43SLOT45END. The projection identity gives ENDPOINTSOURCE0C01A7DD4A656A73579A2F43SLOT46END and ENDPOINTSOURCE0C01A7DD4A656A73579A2F43SLOT47END. Since the ENDPOINTSOURCE0C01A7DD4A656A73579A2F43SLOT48END values ENDPOINTSOURCE0C01A7DD4A656A73579A2F43SLOT49END are positive, the ENDPOINTSOURCE0C01A7DD4A656A73579A2F43SLOT50END form an orthonormal basis of ENDPOINTSOURCE0C01A7DD4A656A73579A2F43SLOT51END. These are principal bases. Set ENDPOINTSOURCE0C01A7DD4A656A73579A2F43SLOT52END. Consequently

ENDPOINTSOURCE0C01A7DD4A656A73579A2F43SLOT53END

The multiplicity of ENDPOINTSOURCE0C01A7DD4A656A73579A2F43SLOT54END is exactly ENDPOINTSOURCE0C01A7DD4A656A73579A2F43SLOT55END. Indeed equality ENDPOINTSOURCE0C01A7DD4A656A73579A2F43SLOT56END for an orthogonal projection is equivalent to ENDPOINTSOURCE0C01A7DD4A656A73579A2F43SLOT57END; apply this on ENDPOINTSOURCE0C01A7DD4A656A73579A2F43SLOT58END. The same multiplicity is ENDPOINTSOURCE0C01A7DD4A656A73579A2F43SLOT59END: a common vector has the same remainder in both sections, and conversely equal section values define a vector in the intersection. Thus all unit cosines are retained, including multiplicities arising from coinciding section subspaces.

Taking the determinant of the actual positive matrix in (PA8) gives

ENDPOINTSOURCE0C01A7DD4A656A73579A2F43SLOT60END

All terms are finite and nonnegative, because the original section Grams are positive definite.

\subsubsection{3. The two pairs and the forced unit cosine}\label{endpoint-principal_angle_volume_review--the-two-pairs-and-the-forced-unit-cosine}

Use exactly the pairs

ENDPOINTSOURCE0C01A7DD4A656A73579A2F43SLOT61END

For the second pair, both section spaces are contained in

ENDPOINTSOURCE0C01A7DD4A656A73579A2F43SLOT62END

For ENDPOINTSOURCE0C01A7DD4A656A73579A2F43SLOT63END, this follows from its defining orthogonality to ENDPOINTSOURCE0C01A7DD4A656A73579A2F43SLOT64END; for ENDPOINTSOURCE0C01A7DD4A656A73579A2F43SLOT65END, its relation space contains ENDPOINTSOURCE0C01A7DD4A656A73579A2F43SLOT66END, so it has the same orthogonality. The dimension in (PA11) is exact because ENDPOINTSOURCE0C01A7DD4A656A73579A2F43SLOT67END is nonzero and its norm is positive. Both section spaces have dimension ENDPOINTSOURCE0C01A7DD4A656A73579A2F43SLOT68END. The identity ENDPOINTSOURCE0C01A7DD4A656A73579A2F43SLOT69END therefore gives an intersection of dimension at least one. By (PA8), at least one angle in the second pair has cosine one. For ENDPOINTSOURCE0C01A7DD4A656A73579A2F43SLOT70END, the second pair is literally ENDPOINTSOURCE0C01A7DD4A656A73579A2F43SLOT71END, so its sole angle is zero and the whole pair contributes zero.

The original four-volume is precisely

ENDPOINTSOURCE0C01A7DD4A656A73579A2F43SLOT72END

Remove one of the proved zero contributions from the second pair. There remain exactly ENDPOINTSOURCE0C01A7DD4A656A73579A2F43SLOT73END nonnegative slots, with every additional zero slot kept. This removal changes neither the value nor either bound below. It is a pointwise assertion; no differentiability of a chosen angle enumeration at eigenvalue crossings is required.

\subsubsection{4. Full projection-difference spectrum and derivative estimate}\label{endpoint-principal_angle_volume_review--full-projection-difference-spectrum-and-derivative-estimate}

For a positive angle in (PA7), define ENDPOINTSOURCE0C01A7DD4A656A73579A2F43SLOT74END. The principal-basis identities show that the ENDPOINTSOURCE0C01A7DD4A656A73579A2F43SLOT75END are orthonormal and orthogonal to ENDPOINTSOURCE0C01A7DD4A656A73579A2F43SLOT76END, and that distinct positive-angle two-planes are orthogonal. In the ordered basis ENDPOINTSOURCE0C01A7DD4A656A73579A2F43SLOT77END, the two projections and their difference have the exact matrices

ENDPOINTSOURCE0C01A7DD4A656A73579A2F43SLOT78END

where ENDPOINTSOURCE0C01A7DD4A656A73579A2F43SLOT79END. The last matrix has trace zero and determinant ENDPOINTSOURCE0C01A7DD4A656A73579A2F43SLOT80END, hence eigenvalues ENDPOINTSOURCE0C01A7DD4A656A73579A2F43SLOT81END. A zero-angle common vector has zero difference, and the orthogonal complement of the two spaces also has zero difference. Thus these are the complete nonzero eigenvalues with multiplicity. In particular

ENDPOINTSOURCE0C01A7DD4A656A73579A2F43SLOT82END

Here the positive and negative parts are positive operators, and the latter is the negative of the operator on its negative spectral subspace.

Let ENDPOINTSOURCE0C01A7DD4A656A73579A2F43SLOT83END, which is self-adjoint for the same ENDPOINTSOURCE0C01A7DD4A656A73579A2F43SLOT84END. Let ENDPOINTSOURCE0C01A7DD4A656A73579A2F43SLOT85END. For any trace-zero self-adjoint ENDPOINTSOURCE0C01A7DD4A656A73579A2F43SLOT86END, subtracting the scalar ENDPOINTSOURCE0C01A7DD4A656A73579A2F43SLOT87END leaves its trace pairing unchanged. The inequalities ENDPOINTSOURCE0C01A7DD4A656A73579A2F43SLOT88END, applied to the two positive parts of ENDPOINTSOURCE0C01A7DD4A656A73579A2F43SLOT89END, give

ENDPOINTSOURCE0C01A7DD4A656A73579A2F43SLOT90END

To verify the inequality without a commutation assumption, the trace of a product of two positive operators is nonnegative: it equals the trace of their positive conjugation ENDPOINTSOURCE0C01A7DD4A656A73579A2F43SLOT91END. Hence each of the two positive-part trace terms lies between zero and ENDPOINTSOURCE0C01A7DD4A656A73579A2F43SLOT92END, and their difference has the asserted absolute bound.

The differentiability of the canonical section follows from its explicit positive-Gram inverse formula. Because ENDPOINTSOURCE0C01A7DD4A656A73579A2F43SLOT93END, its derivative is a relation vector. Orthogonality to that relation space cancels the two section-derivative cross terms, giving

ENDPOINTSOURCE0C01A7DD4A656A73579A2F43SLOT94END

The last formula is the orthogonal projection onto ENDPOINTSOURCE0C01A7DD4A656A73579A2F43SLOT95END, since it is the identity there and annihilates its orthogonal complement. Cyclicity of the finite trace then proves the middle identity.

Using the two pairs (PA10), (PA12), (PA14)--(PA16), and the triangle inequality gives

ENDPOINTSOURCE0C01A7DD4A656A73579A2F43SLOT96END

This bound does not presume that the two projection differences commute, or that either one commutes with ENDPOINTSOURCE0C01A7DD4A656A73579A2F43SLOT97END.

For each of the ENDPOINTSOURCE0C01A7DD4A656A73579A2F43SLOT98END retained slots, set ENDPOINTSOURCE0C01A7DD4A656A73579A2F43SLOT99END. Then ENDPOINTSOURCE0C01A7DD4A656A73579A2F43SLOT100END and ENDPOINTSOURCE0C01A7DD4A656A73579A2F43SLOT101END, where

ENDPOINTSOURCE0C01A7DD4A656A73579A2F43SLOT102END

It is continuous at zero and strictly concave on ENDPOINTSOURCE0C01A7DD4A656A73579A2F43SLOT103END, because

ENDPOINTSOURCE0C01A7DD4A656A73579A2F43SLOT104END

Concavity extends to the endpoint by continuity: apply the interior inequality to all arguments increased by a positive number and let that number decrease to zero. Jensen's inequality, which follows by iterating the defining two-point concavity inequality, therefore gives

ENDPOINTSOURCE0C01A7DD4A656A73579A2F43SLOT105END

Combining this with (PA17) proves the proposed estimate with its exact constant:

ENDPOINTSOURCE0C01A7DD4A656A73579A2F43SLOT106END

\subsubsection{5. Strict positivity for the actual moment Grams}\label{endpoint-principal_angle_volume_review--strict-positivity-for-the-actual-moment-grams}

The original arithmetic and Gamma norms have the common vertical coordinate ENDPOINTSOURCE0C01A7DD4A656A73579A2F43SLOT107END and the form

ENDPOINTSOURCE0C01A7DD4A656A73579A2F43SLOT108END

The original densities and their convex interpolation have the required finite moments and positive definite polynomial Grams. Write the unchanged polynomial as ENDPOINTSOURCE0C01A7DD4A656A73579A2F43SLOT109END, retaining its coefficients, and define the exact reflected polynomial

ENDPOINTSOURCE0C01A7DD4A656A73579A2F43SLOT110END

On the original vertical line, ENDPOINTSOURCE0C01A7DD4A656A73579A2F43SLOT111END, so ENDPOINTSOURCE0C01A7DD4A656A73579A2F43SLOT112END. It has degree ENDPOINTSOURCE0C01A7DD4A656A73579A2F43SLOT113END, with leading coefficient ENDPOINTSOURCE0C01A7DD4A656A73579A2F43SLOT114END; that sign is not removed. Hence ENDPOINTSOURCE0C01A7DD4A656A73579A2F43SLOT115END and

ENDPOINTSOURCE0C01A7DD4A656A73579A2F43SLOT116END

All integrals are within the specified moment degree ENDPOINTSOURCE0C01A7DD4A656A73579A2F43SLOT117END. If the first pair in (PA10) had every cosine equal to one, its two ENDPOINTSOURCE0C01A7DD4A656A73579A2F43SLOT118END-dimensional spaces would coincide. Its first space is exactly ENDPOINTSOURCE0C01A7DD4A656A73579A2F43SLOT119END, and its second is ENDPOINTSOURCE0C01A7DD4A656A73579A2F43SLOT120END. The constant polynomial would then be orthogonal to ENDPOINTSOURCE0C01A7DD4A656A73579A2F43SLOT121END, contradicting (PA22). Thus at least one angle in the first pair is positive, and

ENDPOINTSOURCE0C01A7DD4A656A73579A2F43SLOT122END

for every actual moment Gram in this path. This proves strict positivity without a presumption about the zero locations or a discarded relation primitive. It also verifies the needed endpoint at ENDPOINTSOURCE0C01A7DD4A656A73579A2F43SLOT123END, where the second pair contributes exactly zero.

\subsubsection{6. Integrated arcosh comparison on the actual arithmetic/reference path}\label{endpoint-principal_angle_volume_review--integrated-arcosh-comparison-on-the-actual-arithmeticreference-path}

For ENDPOINTSOURCE0C01A7DD4A656A73579A2F43SLOT124END, define

ENDPOINTSOURCE0C01A7DD4A656A73579A2F43SLOT125END

The derivative is obtained by differentiating the displayed exponential and ENDPOINTSOURCE0C01A7DD4A656A73579A2F43SLOT126END; no rescaling of ENDPOINTSOURCE0C01A7DD4A656A73579A2F43SLOT127END or its original mass is performed. By (PA23), the chain rule applies along the entire actual path. Equation (PA19) gives

ENDPOINTSOURCE0C01A7DD4A656A73579A2F43SLOT128END

Now take exactly the convex source-metric interpolation ENDPOINTSOURCE0C01A7DD4A656A73579A2F43SLOT129END, ENDPOINTSOURCE0C01A7DD4A656A73579A2F43SLOT130END, inherited from the two actual densities in (PA20). Put

ENDPOINTSOURCE0C01A7DD4A656A73579A2F43SLOT131END

This is the complete common-source relative Gram, not a quotient-metric substitute. Factoring ENDPOINTSOURCE0C01A7DD4A656A73579A2F43SLOT132END shows that ENDPOINTSOURCE0C01A7DD4A656A73579A2F43SLOT133END is similar to ENDPOINTSOURCE0C01A7DD4A656A73579A2F43SLOT134END. Its eigenvalues are

ENDPOINTSOURCE0C01A7DD4A656A73579A2F43SLOT135END

Every denominator is positive. For fixed ENDPOINTSOURCE0C01A7DD4A656A73579A2F43SLOT136END, the derivative of this scalar expression with respect to ENDPOINTSOURCE0C01A7DD4A656A73579A2F43SLOT137END is ENDPOINTSOURCE0C01A7DD4A656A73579A2F43SLOT138END, so its largest and smallest values have the same indices as the original ENDPOINTSOURCE0C01A7DD4A656A73579A2F43SLOT139END. Therefore

ENDPOINTSOURCE0C01A7DD4A656A73579A2F43SLOT140END

The identity also includes any ENDPOINTSOURCE0C01A7DD4A656A73579A2F43SLOT141END, whose integrand is zero, and cases in which both extreme eigenvalues are above or below one.

Integrating (PA25) and using (PA28) gives precisely

ENDPOINTSOURCE0C01A7DD4A656A73579A2F43SLOT142END

Every endpoint volume is the original four-volume in (PA12). All section angles are strictly below ENDPOINTSOURCE0C01A7DD4A656A73579A2F43SLOT143END because every finite quotient Gram is positive definite. The forced unit-cosine slot is retained through the exact common-space argument (PA11). For ENDPOINTSOURCE0C01A7DD4A656A73579A2F43SLOT144END, ENDPOINTSOURCE0C01A7DD4A656A73579A2F43SLOT145END and the formula uses the single nontrivial pair. If ENDPOINTSOURCE0C01A7DD4A656A73579A2F43SLOT146END, the two source Grams differ by one positive scalar; all quotient Grams acquire that same scalar and their four-volume ratio is unchanged, consistently giving zero on both sides of (PA29).

The proposed formulas are accepted. This calculation proves their exact finite metric geometry and their application to the original positive moment interpolation. It supplies no bound on ENDPOINTSOURCE0C01A7DD4A656A73579A2F43SLOT147END beyond its computed source definition, and no extra arithmetic conclusion is assigned to that quantity.
